\chapter{Analysis}\label{ch:analysis}\label{ch:2}

Chapter~\ref{ch:req} described the business and fixed the boundary of the system. This chapter analyses what the
system must do and how it behaves, following Lectures~02--04~\cite{slidesL02,slidesL03,slidesL04}. It first turns the
needs into \nFR{} functional and \nNFR{} non-functional requirements, prioritises them and identifies the actors and
\nUseCases{} use cases (Lecture~02). It then answers the central question of Lecture~03: \emph{when an actor
performs a use case, what happens?} -- each main use case is specified and followed directly by the diagram that
shows its behaviour. Finally, following Lecture~04 (Domain Model and Analysis Class Diagram), it moves from
behaviour to structure -- which concepts the system must represent, how they are related and which of them must
be stored (the conceptual and logical data model) -- and checks that
requirements, use cases, scenarios and models trace to each other. The whole chapter stays at the \emph{analysis}
level: there are no controllers, services, repositories or tables yet -- they belong to the design in Chapter~4.

%% =========================================================================
\section{Requirement Levels}\label{sec:1.8}

Lecture~02 starts with the question: what must the new system do? A \emph{requirement} describes a capability,
condition or constraint that the system must satisfy. Requirements have three levels~\cite{slidesL02,sommerville}: \emph{business
requirements} say why the organisation needs the system; \emph{user requirements} say what users need to
achieve; \emph{system requirements} say what the system must provide. Table~\ref{tab:levels} shows how the same
problem moves through the three levels.

\begin{longtable}{@{}P{2.4cm}P{4.0cm}P{4.2cm}P{4.4cm}@{}}
\caption{The three requirement levels of Họa.}\label{tab:levels}\\
\toprule
\textbf{Problem} & \textbf{Business requirement} & \textbf{User requirement} & \textbf{System requirement}\\
\midrule
\endfirsthead
\toprule
\textbf{Problem} & \textbf{Business requirement} & \textbf{User requirement} & \textbf{System requirement}\\
\midrule
\endhead
\small P1 hard to find paintings & \small More customers find a painting and buy it (BO1, BO2) & \small Customers need to search by description, voice or photo, or by a photo of their room & \small FR01, FR02, FR03, FR04, FR14, NFR08\\
\small P2 sold-out prints sold & \small No order cancelled because of stock (BO3) & \small Customers need to know that a print is in stock before buying & \small FR05, FR07, FR12\\
\small P3 order status unknown & \small Fewer order questions (BO4) & \small Customers need to see the status of their orders at any time & \small FR08, FR09, FR10, FR13\\
\small P4 manual work & \small No wrong prices, less staff time (BO3, BO4) & \small The administrator needs to manage paintings and order status in one place & \small FR12, FR13, FR20\\
\small P5 risk of data leaks & \small Protect customer data (BO5) & \small Customers need to be sure that nobody else can see their orders & \small NFR01, NFR05\\
\bottomrule
\end{longtable}

\subsection{Stakeholder needs}\label{sec:needs}

Assignment~02 puts one more step between a problem and a requirement: the \emph{stakeholder need}. A need says
what a stakeholder must be able to do or to know; it is not yet a solution (``find a painting without knowing its
title'', not ``use Gemma~4''). Table~\ref{tab:needs} derives the needs from the stakeholder analysis
(Table~\ref{tab:stakeholders}) and the problem analysis (Table~\ref{tab:pcin}) of Chapter~\ref{ch:req}. The third
column gives the problem behind each need, or the reason when the need does not come from one of the problems
P1--P5; the last column names the requirements of Section~\ref{sec:frnfr} that answer it. Each of the \nFR{}
functional and \nNFR{} non-functional requirements appears at least once, so no requirement exists without a
stakeholder who needs it.

\begin{longtable}{@{}P{2.5cm}P{4.4cm}P{5.3cm}P{2.8cm}@{}}
\caption{Stakeholder needs and the requirements that answer them}\label{tab:needs}\\
\toprule
\textbf{Stakeholder} & \textbf{Need} & \textbf{Reason or problem} & \textbf{Requirements}\\
\midrule
\endfirsthead
\toprule
\textbf{Stakeholder} & \textbf{Need} & \textbf{Reason or problem} & \textbf{Requirements}\\
\midrule
\endhead
\small Customer & \small Find a painting without knowing its title: from a description, by voice or with a photo & \small P1: the catalogue can be searched by title only, while customers describe a colour or a subject, or send a photo & \small FR01--FR04, NFR07\\
\small Customer & \small Choose a painting that suits their own room & \small P1: customers know the room they decorate better than the painting that would fit it & \small FR14\\
\small Customer & \small Understand why a painting is suggested & \small P1: without a salesperson to explain it, a suggestion is hard to trust and the consultation gets long & \small FR18\\
\small Customer & \small Get a useful answer from every action, even when the AI service is slow or down & \small Conflict of Section~\ref{sec:1.3}: search at any time, but the AI services can fail; P3: an unclear answer becomes one more question to staff & \small NFR02, NFR09, NFR12\\
\small Customer & \small Collect several prints and frames before paying & \small P4: today the order is put together in messages with a member of staff & \small FR06\\
\small Customer, warehouse staff & \small Buy, and pack, only what is really in stock & \small P2: stock is updated by hand, apart from the orders, so sold-out prints are sold & \small FR05, FR07, NFR11\\
\small Customer, shop owner & \small One correct total for every order, with the VIP discount and the delivery cost & \small P4: prices and totals are worked out by hand, so wrong prices occur & \small FR15, FR17\\
\small Customer & \small Give the delivery address while ordering & \small P4: order details, the address included, are copied by hand from messages into a spreadsheet & \small FR16\\
\small Customer & \small See the status of their orders at any time without asking staff & \small P3: the status exists only in a staff spreadsheet & \small FR08, FR09, FR10\\
\small Customer & \small Cancel an order that has not left the shop & \small P4: orders, cancellations included, are handled by staff by hand & \small FR11\\
\small Customer & \small Be sure that nobody else sees their orders, card or voice & \small P5: customers are identified by an order id in a message; card numbers are written down & \small NFR01, NFR05, NFR08\\
\small Support staff & \small Spend less time answering ``where is my order?'' & \small P3: repeated questions overload the staff & \small FR08, FR10\\
\small Shop administrator & \small Keep paintings, prices and stock right in one place & \small P2, P4: catalogue and stock are kept in separate spreadsheets updated by hand & \small FR12\\
\small Shop administrator & \small Accept only customer photos that the system can read safely & \small Photo search (P1) brings files from outside the shop into it; a wrong or very large file must not break the search & \small NFR03\\
\small Shop administrator, warehouse staff & \small See the orders of all customers and move each one to its next status & \small P3, P4: orders are processed by hand and their status is known only to staff & \small FR13, FR20\\
\small Shop owner & \small Manage prints and frames as two product lines, each with its own details & \small A print has an artist, a style and a size, a frame a material and a colour; one list would mix them & \small FR19, NFR11\\
\small Developer, lecturer & \small Models and code that say the same thing, and a demo that runs on one laptop & \small The essay is checked against the lectures; the demo must run with the network needed only for Gemma & \small NFR04, NFR06, NFR10\\
\bottomrule
\end{longtable}

%% =========================================================================
\section{Functional and Non-functional Requirements}\label{sec:frnfr}\label{sec:1.9}

A \emph{functional requirement} describes a function or behaviour that the system provides. A
\emph{non-functional requirement} describes a quality, constraint or condition of the system: performance,
security, reliability, maintainability. Each requirement follows the pattern of Lecture~02:
\textbf{``The system shall + action + object + condition''}, with an identifier, an area, a priority and a source
(the stakeholder who asks for it). Table~\ref{tab:fr} lists the \nFR{} functional and Table~\ref{tab:nfr} the
\nNFR{} non-functional requirements. The Area column plays the role of the requirement category of Assignment~02: a
functional area (search, catalogue, cart and checkout, orders, assistant, administration) for a functional
requirement and a quality (security, reliability, privacy, performance, usability, maintainability, portability)
for a non-functional one; it is not the product \code{Category} of FR19. FR15--FR20 and NFR09--NFR10 add no new
feature but make explicit what the shop already did: the VIP discount, the address chosen at checkout, the shipping
fee, the reason shown with each suggestion, the product categories and the list of all orders for the
administrator, and the usability and maintainability rules the code already followed. The two tables are generated
from the shared specification file of the project, so the identifiers are the same as in the diagrams, the tests
and the code.

\begin{longtable}{@{}P{1.2cm}P{6.7cm}P{0.9cm}P{1.3cm}P{2.6cm}P{1.8cm}@{}}
\caption{Functional requirements (ID / Type / Priority / Source)}\label{tab:fr}\\
\toprule
\small \textbf{ID} & \small \textbf{Requirement} & \small \textbf{Type} & \small \textbf{Priority} & \small \textbf{Source} & \small \textbf{Use case}\\
\midrule
\endfirsthead
\toprule
\small \textbf{ID} & \small \textbf{Requirement} & \small \textbf{Type} & \small \textbf{Priority} & \small \textbf{Source} & \small \textbf{Use case}\\
\midrule
\endhead
\bottomrule
\endlastfoot
\small \textbf{FR01} & \small The system shall allow customers to search paintings by keyword or natural sentence and filter them by style, subject, colour, artist and maximum price. & \small FR & \small Must & \small Customer & \small UC02\\
\small \textbf{FR02} & \small The system shall accept a spoken search query, transcribe it with the self-hosted Whisper model and let the customer edit the transcript before searching. & \small FR & \small Should & \small Customer & \small UC03\\
\small \textbf{FR03} & \small The system shall return the paintings most similar to a photo uploaded by the customer, ranked by similarity. & \small FR & \small Should & \small Customer & \small UC04\\
\small \textbf{FR04} & \small The system shall combine a text query and an image in one search. & \small FR & \small Could & \small Customer & \small UC01\\
\small \textbf{FR05} & \small The system shall display painting details including artist, year, size, medium, price and current stock. & \small FR & \small Must & \small Customer & \small UC05\\
\small \textbf{FR06} & \small The system shall allow customers to add paintings to the cart, change quantities and remove cart lines. & \small FR & \small Must & \small Customer & \small UC06\\
\small \textbf{FR07} & \small The system shall create an order from the cart when the paintings are in stock, freeze the prices, decrease the stock and empty the cart. & \small FR & \small Must & \small Customer, Warehouse & \small UC07\\
\small \textbf{FR08} & \small The system shall allow customers to view their orders and the current status of each order. & \small FR & \small Must & \small Customer & \small UC08\\
\small \textbf{FR09} & \small The system shall allow customers to look up an order by id only when the order belongs to their account. & \small FR & \small Must & \small Customer & \small UC09\\
\small \textbf{FR10} & \small The system shall answer shopping questions with an AI assistant that can call the painting search and order look-up functions. & \small FR & \small Should & \small Customer, Support staff & \small UC10\\
\small \textbf{FR11} & \small The system shall allow customers to cancel an order before it is shipped. & \small FR & \small Should & \small Customer & \small UC11\\
\small \textbf{FR12} & \small The system shall allow the administrator to add paintings and to update their price and stock. & \small FR & \small Must & \small Administrator & \small UC12\\
\small \textbf{FR13} & \small The system shall allow the administrator to update the processing and delivery status of orders. & \small FR & \small Must & \small Administrator, Warehouse & \small UC13\\
\small \textbf{FR14} & \small The system shall suggest paintings that suit a room from a photo of the room and let the customer preview a painting on the wall. & \small FR & \small Could & \small Customer & \small UC14\\
\end{longtable}


\begin{longtable}{@{}P{1.2cm}P{5.5cm}P{2.4cm}P{1.35cm}P{2.3cm}P{1.75cm}@{}}
\caption{Non-functional requirements (ID / Area / Priority / Source)}\label{tab:nfr}\\
\toprule
\small \textbf{ID} & \small \textbf{Requirement} & \small \textbf{Area} & \small \textbf{Priority} & \small \textbf{Source} & \small \textbf{Use case}\\
\midrule
\endfirsthead
\toprule
\small \textbf{ID} & \small \textbf{Requirement} & \small \textbf{Area} & \small \textbf{Priority} & \small \textbf{Source} & \small \textbf{Use case}\\
\midrule
\endhead
\bottomrule
\endlastfoot
\small \textbf{NFR01} & \small The system shall take the identity from the signed-in session and check order ownership on the server. & \small Security & \small Must & \small Customer, Administrator & \small UC08, UC09, UC10\\
\small \textbf{NFR02} & \small The system shall still find paintings with keyword rules when the Gemma service fails or times out. & \small Reliability & \small Must & \small Customer & \small UC02, UC03, UC04, UC10\\
\small \textbf{NFR03} & \small The system shall reject uploaded images that are not JPEG, PNG or WebP or are larger than 5 MB. & \small Security & \small Must & \small Administrator & \small UC04\\
\small \textbf{NFR04} & \small The system shall separate presentation, application, domain and infrastructure; the user interface shall not access the database directly. & \small Maintainability & \small Must & \small Developer & \small UC01\\
\small \textbf{NFR05} & \small The system shall not store full card numbers; only the last four digits and the authorisation code are kept. & \small Privacy & \small Must & \small Customer & \small UC07\\
\small \textbf{NFR06} & \small The system shall run on one machine with a SQLite database and need the network only to call Gemma. & \small Portability & \small Should & \small Lecturer, Developer & \small UC01\\
\small \textbf{NFR07} & \small The system shall return text search results within 2 seconds on the demo data. & \small Performance & \small Should & \small Customer & \small UC02\\
\small \textbf{NFR08} & \small The system shall process audio on the local machine with self-hosted Whisper, never send audio outside, and transcribe a 5-second sentence within 3 seconds. & \small Privacy & \small Must & \small Customer & \small UC03\\
\small \textbf{NFR09} & \small The system shall end every search, cart and order action in a visible result, an empty state with a hint, or an error message that says what to do. & \small Usability & \small Should & \small Customer & \small UC02, UC06, UC07\\
\small \textbf{NFR10} & \small The system shall keep the class, method, table and column names of the code identical to the design model and check this automatically before each release. & \small Maintainability & \small Must & \small Developer, Lecturer & \small UC01\\
\small \textbf{NFR11} & \small The system shall enforce every key and constraint of the physical data model, never let stock become negative, and keep the details of each product in the table its category names. & \small Data integrity & \small Must & \small Shop owner, Administrator & \small UC07, UC12\\
\small \textbf{NFR12} & \small The system shall keep the catalogue, cart, checkout, orders and administration working when the Gemma or Whisper service is unavailable. & \small Availability & \small Must & \small Customer, Shop owner & \small UC05, UC07, UC08, UC13\\
\end{longtable}


A common mistake pointed out in Lecture~02 is to write ``the system shall be fast''. How many seconds, under which
load? NFR07 therefore gives a limit and a condition (first row of Table~\ref{tab:revised}), and the
waiting limit of an AI call is set separately, because the response time of an external service is not under the
control of Họa. NFR08 is both a privacy requirement (audio never leaves the machine) and a measurable time limit,
so it decides how speech recognition is built in Chapter~3. Table~\ref{tab:nfrcheck} in Appendix~C states how
each non-functional requirement is verified.

\subsection{Criteria of a good requirement}

A good requirement is correct, complete, consistent, unambiguous, feasible, verifiable, traceable, prioritised,
understandable and testable~\cite{iso29148,slidesL02}. The simplest test is the question of Lecture~02: \emph{can this requirement be
verified objectively?} Table~\ref{tab:goodreq} checks the criteria against the requirements of Họa.

\begin{longtable}{@{}P{2.6cm}P{12.4cm}@{}}
\caption{Checking the criteria of a good requirement}\label{tab:goodreq}\\
\toprule
\textbf{Criterion} & \textbf{How the requirements of Họa. meet it}\\
\midrule
\endfirsthead
\toprule
\textbf{Criterion} & \textbf{How the requirements of Họa. meet it}\\
\midrule
\endhead
\small Unambiguous & \small Each statement has one main action; FR07 lists the four things that must happen at checkout instead of saying ``process the order''.\\
\small Verifiable & \small NFR07 has a 2-second limit and the condition ``on the demo data''; NFR03 names the formats and the size; NFR08 names the place of processing and a 3-second limit.\\
\small Consistent & \small FR11 (cancel before shipment) agrees with the order statuses of FR13 and with the state machine of Section~\ref{sec:2.9}.\\
\small Feasible & \small Image search uses an existing model through an API instead of training one; speech recognition uses an open model that runs on a laptop.\\
\small Traceable & \small Every requirement has a stable id and a list of use cases; the matrix is in Section~\ref{sec:trace}.\\
\small Prioritised & \small Every requirement is Must, Should or Could (Table~\ref{tab:moscow}).\\
\small Testable & \small The Given -- When -- Then acceptance criteria of Table~\ref{tab:stories} become the test and demo scenarios of Chapter~\ref{ch:code}.\\
\bottomrule
\end{longtable}

Assignment~02 also asks for the opposite check: find requirements that break these criteria and revise them. The
requirements of Họa. did not start in their present form. Table~\ref{tab:revised} records five weak statements
met while the project grew, the defect of each, and the requirement as it now stands in the shared specification.

\begin{longtable}{@{}P{3.9cm}P{4.4cm}P{6.7cm}@{}}
\caption{Requirement validation: weak statements and their revision}\label{tab:revised}\\
\toprule
\textbf{Weak statement} & \textbf{Problem} & \textbf{Revised requirement (current model)}\\
\midrule
\endfirsthead
\toprule
\textbf{Weak statement} & \textbf{Problem} & \textbf{Revised requirement (current model)}\\
\midrule
\endhead
\small ``The search shall be fast.'' & \small Not measurable: no time limit and no condition, so no test can pass or fail it. & \small NFR07: The system shall return text search results within 2 seconds on the demo data.\\
\small ``The system shall use Gemma~4 to search paintings.'' & \small A solution, not a requirement: it names a technology instead of what the customer needs, and says nothing when the model fails. & \small FR01: The system shall allow customers to search paintings by keyword or natural sentence and filter them by style, subject, colour, artist and maximum price. FR03: The system shall return the paintings most similar to a photo uploaded by the customer, ranked by similarity. Gemma~4 is a design choice (Chapter~3); its failure is covered by NFR02.\\
\small ``Find an order by its id'', the starting point taken from \code{find_order(order_id)} of Assignment~06. & \small Incomplete and unsafe: no access condition, so anyone who knows or guesses an id sees another customer's order (P5). & \small FR09: The system shall allow customers to look up an order by id only when the order belongs to their account. NFR01 adds where the identity comes from: the signed-in session, with ownership checked on the server.\\
\small ``Search by voice'' through a cloud speech service, the first plan (change CR-04). & \small Incomplete: it says nothing about where the audio goes, although a recording of a voice is personal data (Section~\ref{sec:1.3}). & \small FR02 keeps the function; NFR08: The system shall process audio on the local machine with self-hosted Whisper, never send audio outside, and transcribe a 5-second sentence within 3 seconds.\\
\small FR07 as the only checkout requirement: create the order, freeze prices, decrease stock, empty the cart. & \small Incomplete: the VIP discount, the choice of address and the shipping fee were done by the shop but written in no requirement, so they could not be traced or checked against one. & \small FR15: The system shall apply the discount rate of a VIP customer to the order total at checkout. FR16: The system shall let the customer choose one of their addresses or add a new address during checkout. FR17: The system shall add a flat shipping fee of 8 USD to every order.\\
\bottomrule
\end{longtable}

\subsection{Acceptance criteria}

An acceptance criterion is the check that tells whether one requirement is met. Table~\ref{tab:stories} writes
the main requirements as user stories with Given -- When -- Then conditions; the queries are the demo queries of
the shared specification, which are the painting-shop versions of the example queries of Assignment~06.

\begin{longtable}{@{}P{1.1cm}P{5.2cm}P{8.7cm}@{}}
\caption{User stories and acceptance criteria}\label{tab:stories}\\
\toprule
\textbf{ID} & \textbf{User story} & \textbf{Acceptance criteria (Given -- When -- Then)}\\
\midrule
\endfirsthead
\toprule
\textbf{ID} & \textbf{User story} & \textbf{Acceptance criteria (Given -- When -- Then)}\\
\midrule
\endhead
\small US01 & \small As a customer, I want to type ``blue landscape paintings'' and see suitable paintings without knowing any title (FR01). & \small Given the catalogue has blue landscapes, When the customer sends the sentence, Then the top results have subject landscape or dominant colour blue, ordered by decreasing score.\\
\small US02 & \small As a customer, I want to search ``Monet paintings under 100 dollars'' (FR01). & \small Given Monet prints at several prices, When the customer sends the sentence, Then every result is a Monet print priced at most 100 USD.\\
\small US03 & \small As a customer, I want to speak instead of typing (FR02, NFR08). & \small Given the browser may use the microphone, When the customer says the query, Then the transcript appears in the search box for editing, the results are the same as when typing the sentence, and no audio is sent outside the machine.\\
\small US04 & \small As a customer, I want to upload a photo and find paintings in a similar style (FR03). & \small Given a valid JPEG under 5 MB, When the customer uploads it, Then the system shows the paintings with the closest style, subject and colour, each with a similarity score.\\
\small US05 & \small As a customer, I want to add paintings to the cart and place an order (FR06, FR07). & \small Given a cart with at least one line and enough stock, When the customer confirms address and payment method, Then a PLACED order is created with frozen prices, stock decreases and the cart is emptied; an empty cart is refused.\\
\small US06 & \small As a customer, I want to look up an order by its id (FR09, NFR01). & \small Given order 20261001 belongs to C001, When C001 asks ``find my order 20261001'', Then the status is shown; When C002 asks for the same id, Then the system answers ``order not found''.\\
\small US07 & \small As a customer, I want to ask the assistant ``where is my latest order?'' (FR10). & \small Given Gemma is available, When the customer asks, Then the assistant calls the order look-up tool for the customer of the session and states the status; if Gemma fails, the page says that the assistant is unavailable.\\
\small US08 & \small As a customer, I want to cancel an order I have just placed (FR11). & \small Given an order in status PLACED, When the customer cancels it, Then it becomes CANCELLED; Given an order already SHIPPED, Then the cancellation is refused.\\
\small US09 & \small As the administrator, I want to change the price and stock of a painting (FR12). & \small Given an administrator session, When a valid price and stock are entered, Then the catalogue shows the new price; a negative price or stock is refused.\\
\small US10 & \small As the administrator, I want to move an order to its next status (FR13). & \small Given an order PLACED, When it is moved to PROCESSING, Then it succeeds; When it is moved directly to DELIVERED, Then it is refused because the order of statuses is wrong.\\
\small US11 & \small As a customer, I want to upload a photo of my living room and see which paintings suit it (FR14). & \small Given a valid room photo, When the customer uploads it, Then the system shows the palette and style it read from the room and suggests three to five paintings, each with a reason, and one can be previewed on the wall of the photo.\\
\bottomrule
\end{longtable}

%% =========================================================================
\section{Prioritization and Change Management}\label{sec:1.11}

Not all requirements are equally important. MoSCoW divides them into Must (essential), Should (important but not
essential), Could (useful improvement) and Won't (not in this release). Besides business value, I
considered risk and dependencies: image search is attractive but depends on an external service, so it is
Should and comes with the fallback requirement NFR02 at level Must. Voice privacy (NFR08) is Must even though
voice search (FR02) is only Should: if voice search is offered at all, it must be private.

\begin{longtable}{@{}P{2.4cm}P{6.3cm}P{6.3cm}@{}}
\caption{MoSCoW prioritisation of the demo release}\label{tab:moscow}\\
\toprule
\small \textbf{Level} & \small \textbf{Requirements} & \small \textbf{Meaning}\\
\midrule
\endfirsthead
\toprule
\small \textbf{Level} & \small \textbf{Requirements} & \small \textbf{Meaning}\\
\midrule
\endhead
\bottomrule
\endlastfoot
\small Must & \small FR01, FR05, FR06, FR07, FR08, FR09, FR12, FR13, FR15, FR17, FR19, FR20, NFR01, NFR02, NFR03, NFR04, NFR05, NFR08, NFR10, NFR11, NFR12 & \small Without it the gallery cannot operate or is not safe.\\
\small Should & \small FR02, FR03, FR10, FR11, FR16, FR18, NFR06, NFR07, NFR09 & \small Important differentiator; has a fallback when missing.\\
\small Could & \small FR04, FR14 & \small Improves the experience, never blocks a purchase.\\
\small Won't (this release) & \small Real payments, real carriers, personalised recommendations, account registration & \small Outside the scope of the essay; recorded for later.\\
\end{longtable}


Requirements do not stay fixed once they are written. While preparing the essay I met five real changes,
recorded as change -- reason -- affected parts (Table~\ref{tab:changes}); each change was followed along the trace
to use cases, classes, tables, code and tests.

\begin{longtable}{@{}P{1.3cm}P{3.9cm}P{4.4cm}P{5.4cm}@{}}
\caption{Requirement change log}\label{tab:changes}\\
\toprule
\textbf{ID} & \textbf{Change} & \textbf{Reason} & \textbf{Affected parts}\\
\midrule
\endfirsthead
\toprule
\textbf{ID} & \textbf{Change} & \textbf{Reason} & \textbf{Affected parts}\\
\midrule
\endhead
\small CR-01 & \small Product domain changed from books, clothes, shoes and electronics to paintings and frames & \small Search by photo and by feeling is most valuable for paintings & \small Product classes, database tables, sample data, demo queries\\
\small CR-02 & \small AI shopping assistant with tool calling added (FR10) & \small Gemma 4 supports tool calling; fewer order questions (BO4) & \small UC10, class \code{AssistantService}, ownership tests\\
\small CR-03 & \small Storage layer renamed from DAO to Repository & \small Match the class names of the A05 slides and the A06 sample code & \small Layered class diagram, source code, Section~4.4 and Chapter~\ref{ch:code}\\
\small CR-04 & \small Voice search through a self-hosted Whisper server instead of a cloud speech service (FR02, NFR08) & \small Customer voice is personal data; works offline & \small Context view, \code{SpeechService}, \code{WhisperClient}, deployment\\
\small CR-05 & \small Product, data and essay in English; prices in USD; ``find a painting for my room'' added (FR14, UC14) & \small Match the A06 queries (``under 100 dollars''); a new reason to buy & \small Demo queries, catalogue, Whisper language, use case diagram, \code{RoomPreview}\\
\bottomrule
\end{longtable}

%% =========================================================================
\section{Actors and Use Cases}\label{sec:1.10}

\subsection{Stakeholder is not actor}

An actor is \emph{an external role that interacts directly with the system}; it is not necessarily a specific
person. Stakeholders and actors are related but not the same: a stakeholder may only be interested in or
affected by the system without using it, while an actor exchanges information with the system directly
(Table~\ref{tab:stakeactor}).

\begin{longtable}{@{}P{3.6cm}P{2.4cm}P{2.8cm}P{6.2cm}@{}}
\caption{Stakeholders and actors of Họa.}\label{tab:stakeactor}\\
\toprule
\textbf{Party} & \textbf{Stakeholder} & \textbf{Actor} & \textbf{Explanation}\\
\midrule
\endfirsthead
\toprule
\textbf{Party} & \textbf{Stakeholder} & \textbf{Actor} & \textbf{Explanation}\\
\midrule
\endhead
\small Customer & \small Yes & \small Customer & \small Both interested and a direct user\\
\small Shop administrator & \small Yes & \small Administrator & \small Manages paintings and orders directly\\
\small Shop owner & \small Yes & \small No & \small Takes decisions and reads reports through the administrator; no direct use in the demo\\
\small Warehouse and support staff & \small Yes & \small No & \small Use order information through the administrator role\\
\small Payment service & \small Yes & \small Payment Service & \small External system that exchanges payment data\\
\small Delivery carrier & \small Yes & \small No & \small Takes the parcel from the warehouse but exchanges no data with the system in this release (Section~\ref{sec:1.7})\\
\small Gemma 4 server & \small Yes & \small Gemma 4 Service & \small External system called to understand sentences, describe images and choose tools\\
\small Whisper speech server & \small No & \small No & \small Part of the system itself (inside the boundary), not an external role\\
\bottomrule
\end{longtable}

\begin{longtable}{@{}P{3.0cm}P{2.4cm}P{6.0cm}P{3.6cm}@{}}
\caption{Actor catalogue}\label{tab:actors}\\
\toprule
\small \textbf{Actor} & \small \textbf{Kind} & \small \textbf{Role} & \small \textbf{Use cases}\\
\midrule
\endfirsthead
\toprule
\small \textbf{Actor} & \small \textbf{Kind} & \small \textbf{Role} & \small \textbf{Use cases}\\
\midrule
\endhead
\bottomrule
\endlastfoot
\small Customer & \small Person & \small Buys paintings: searches, views, orders, tracks and cancels own orders & \small UC01, UC02, UC03, UC04, UC05, UC06, UC07, UC08, UC09, UC10, UC11, UC14\\
\small Administrator & \small Person & \small Runs the gallery: manages paintings, price, stock and order status & \small UC12, UC13\\
\small Gemma 4 Service & \small External system & \small Understands sentences, describes photos, chooses assistant tools & \small UC02, UC03, UC04, UC10, UC14\\
\small Payment Service & \small External system & \small Confirms card payments and returns an authorisation code (simulated) & \small UC07\\
\end{longtable}


\subsection{A use case is a goal, not a button}

A use case describes a goal-oriented interaction between an actor and the system~\cite{larman}. The key question is:
\emph{what goal does the actor want to achieve?} ``Click the search button'' is not a use case; the goal is
``find a painting''. Table~\ref{tab:notbutton} shows how I grouped user-interface actions into goals.

\begin{longtable}{@{}P{6.6cm}P{8.4cm}@{}}
\caption{User-interface actions and the use cases they belong to}\label{tab:notbutton}\\
\toprule
\textbf{User-interface action (not a use case)} & \textbf{Use case (goal of the actor)}\\
\midrule
\endfirsthead
\toprule
\textbf{User-interface action (not a use case)} & \textbf{Use case (goal of the actor)}\\
\midrule
\endhead
\small Type in the search box and press Enter & \small UC02 Search by Keyword\\
\small Press the microphone button, speak, correct the transcript & \small UC03 Search by Voice\\
\small Choose an image file, look at the preview, press search & \small UC04 Search by Image\\
\small Press ``Add to cart'', change a quantity, remove a line & \small UC06 Manage Cart\\
\small Choose an address and a payment method, press ``Place order'' & \small UC07 Checkout\\
\small Type or speak a question in the curator panel & \small UC10 Ask Shopping Assistant\\
\small Press ``Cancel order'' on the order page & \small UC11 Cancel Order\\
\small Edit the price and stock cells in the admin table & \small UC12 Manage Products\\
\small Upload a room photo, pick a painting from the strip, press ``Try it in my room'' on a painting page & \small UC14 Find Painting for My Room\\
\bottomrule
\end{longtable}

\subsection{Use case catalogue and use case diagram}

The use case names follow the Visual Paradigm diagram and the wording of Assignment~06 (Search Product, Search by
Keyword, Search by Voice, Search by Image, Search Order, View Product, View Order), extended with the use cases of
the purchase flow and of administration. The priority of a use case is the priority of the functional
requirements it realises.

\begin{longtable}{@{}P{1.0cm}P{3.0cm}P{3.8cm}P{2.2cm}P{2.0cm}P{1.2cm}P{1.3cm}@{}}
\caption{Use case catalogue}\label{tab:usecases}\\
\toprule
\footnotesize \textbf{ID} & \footnotesize \textbf{Use case} & \footnotesize \textbf{Goal} & \footnotesize \textbf{Primary actor} & \footnotesize \textbf{Supporting actor} & \footnotesize \textbf{Priority} & \footnotesize \textbf{Req.}\\
\midrule
\endfirsthead
\toprule
\footnotesize \textbf{ID} & \footnotesize \textbf{Use case} & \footnotesize \textbf{Goal} & \footnotesize \textbf{Primary actor} & \footnotesize \textbf{Supporting actor} & \footnotesize \textbf{Priority} & \footnotesize \textbf{Req.}\\
\midrule
\endhead
\bottomrule
\endlastfoot
\footnotesize \textbf{UC01} & \footnotesize Search Product \emph{(abstract)} & \footnotesize Find paintings (general use case) & \footnotesize Customer & \footnotesize -- & \footnotesize Must & \footnotesize FR04, NFR04, NFR06, FR19, NFR10\\
\footnotesize \textbf{UC02} & \footnotesize Search by Keyword & \footnotesize Find paintings with keywords or a natural sentence & \footnotesize Customer & \footnotesize Gemma 4 Service & \footnotesize Must & \footnotesize FR01, NFR02, NFR07, FR18, NFR09\\
\footnotesize \textbf{UC03} & \footnotesize Search by Voice & \footnotesize Find paintings by speaking & \footnotesize Customer & \footnotesize Gemma 4 Service & \footnotesize Should & \footnotesize FR02, NFR02, NFR08\\
\footnotesize \textbf{UC04} & \footnotesize Search by Image & \footnotesize Find paintings that look like a photo & \footnotesize Customer & \footnotesize Gemma 4 Service & \footnotesize Should & \footnotesize FR03, NFR02, NFR03, FR18\\
\footnotesize \textbf{UC05} & \footnotesize View Product & \footnotesize Decide on one painting & \footnotesize Customer & \footnotesize -- & \footnotesize Must & \footnotesize FR05, FR19, NFR12\\
\footnotesize \textbf{UC06} & \footnotesize Manage Cart & \footnotesize Prepare the purchase & \footnotesize Customer & \footnotesize -- & \footnotesize Must & \footnotesize FR06, NFR09\\
\footnotesize \textbf{UC07} & \footnotesize Checkout & \footnotesize Turn the cart into an order & \footnotesize Customer & \footnotesize Payment Service & \footnotesize Must & \footnotesize FR07, NFR05, FR15, FR16, FR17, NFR09, NFR11, NFR12\\
\footnotesize \textbf{UC08} & \footnotesize View Order & \footnotesize Follow my orders & \footnotesize Customer & \footnotesize -- & \footnotesize Must & \footnotesize FR08, NFR01, NFR12\\
\footnotesize \textbf{UC09} & \footnotesize Search Order & \footnotesize Check one of my orders by its id & \footnotesize Customer & \footnotesize -- & \footnotesize Must & \footnotesize FR09, NFR01\\
\footnotesize \textbf{UC10} & \footnotesize Ask Shopping Assistant & \footnotesize Get help from the AI curator & \footnotesize Customer & \footnotesize Gemma 4 Service & \footnotesize Should & \footnotesize FR10, NFR01, NFR02\\
\footnotesize \textbf{UC11} & \footnotesize Cancel Order & \footnotesize Stop an order that is not shipped & \footnotesize Customer & \footnotesize -- & \footnotesize Should & \footnotesize FR11\\
\footnotesize \textbf{UC12} & \footnotesize Manage Products & \footnotesize Keep the catalogue, price and stock right & \footnotesize Administrator & \footnotesize -- & \footnotesize Must & \footnotesize FR12, FR19, NFR11\\
\footnotesize \textbf{UC13} & \footnotesize Update Order Status & \footnotesize Move an order to its next status & \footnotesize Administrator & \footnotesize -- & \footnotesize Must & \footnotesize FR13, FR20, NFR12\\
\footnotesize \textbf{UC14} & \footnotesize Find Painting for My Room & \footnotesize Get paintings that suit my room & \footnotesize Customer & \footnotesize Gemma 4 Service & \footnotesize Should & \footnotesize FR14, FR18\\
\end{longtable}


\landscapefig{c1_usecase}{Use case diagram of Họa.}{Use case diagram of Họa.: four actors and the use cases inside the Họa. Gallery boundary. Visual Paradigm diagram \vpsrc{c1_usecase}.}{\label{fig:c1_usecase}\label{fig:usecase}}

Figure~\ref{fig:usecase} uses the three kinds of UML relationship~\cite{uml}; Table~\ref{tab:uclinks} in Appendix~C
lists every relationship with its meaning.
\emph{Generalization} says that Search by Keyword, Search by Voice and Search by Image are three concrete ways of
the general goal Search Product -- exactly the structure of Assignment~06; Find Painting for My Room is a fourth
way whose input is a photo of a room instead of a photo of a painting. \emph{Include} says that a behaviour is
always used: Checkout always reads and checks the cart (Manage Cart); the assistant always goes through the
painting search or the order look-up tool. \emph{Extend} says that a behaviour is optional: from the search
results the customer may open a painting; while viewing an order that is not shipped, the customer may cancel it.
Gemma 4 Service and Payment Service are supporting actors and take part only in the use cases that need them.

The overview is detailed in three diagrams, one per area of the shop. Figure~\ref{fig:c2_uc_search} enlarges the
search part: UC01 Search Product is abstract (in italics) because a customer always searches in one concrete way --
by keyword, by voice, by image or by a photo of the room -- and these four use cases specialise it; opening a
painting (UC05) is an optional extension of a result list, and Gemma~4 supports the four concrete searches.
Figure~\ref{fig:c2_uc_purchase} details the buying part: checkout always works on the cart (include), and
cancelling is an optional step started from the order page (extend); the payment service authorises card payments.
Figure~\ref{fig:c2_uc_assistant_admin} covers the assistant, which answers a question by running the painting search
and the order look-up for the customer, and the administrator, who keeps the catalogue right and moves orders
through their statuses.

\begin{figure}[H]
\centering
\vpfig{c2_uc_search}{0.30\textheight}
\caption[Use cases: search and discovery]{Use cases of search and discovery: the abstract Search Product with its four concrete searches, View Product as an extension, and Gemma 4 Service supporting the AI-based searches. Visual Paradigm diagram \vpsrc{c2_uc_search}.}\label{fig:c2_uc_search}
\end{figure}
\begin{figure}[H]
\centering
\vpfig{c2_uc_purchase}{0.34\textheight}
\caption[Use cases: cart, checkout and orders]{Use cases of the cart, checkout and orders, with Payment Service supporting checkout. Visual Paradigm diagram \vpsrc{c2_uc_purchase}.}\label{fig:c2_uc_purchase}
\end{figure}
\begin{figure}[H]
\centering
\vpfig{c2_uc_assistant_admin}{0.34\textheight}
\caption[Use cases: assistant and administration]{Use cases of the shopping assistant (it includes the painting search and the order look-up and is supported by Gemma 4 Service) and of the administrator. Visual Paradigm diagram \vpsrc{c2_uc_assistant_admin}.}\label{fig:c2_uc_assistant_admin}
\end{figure}


%% =========================================================================
\section{From Requirement to Behavior}\label{sec:2.1}

A requirement says what is needed but not the whole behaviour. FR07 says ``the system shall create an order from
the cart when the paintings are in stock, freeze the prices, decrease the stock and empty the cart'', but it does
not say what the customer does first, what the system checks, or what happens when the card is refused.
Functional analysis turns requirements into behaviour along the chain of Lecture~03:
\begin{center}
\emph{Requirement (FR) $\rightarrow$ Use case (UC) $\rightarrow$ Scenario $\rightarrow$ Activity model (ACT).}
\end{center}
A good analysis answers four questions: what is the goal of the actor; what is the normal flow; what can go
wrong; how does the system respond. Lecture~03 also lists the usual mistakes: drawing diagrams before understanding
the process, modelling every mouse click, ignoring alternative and error flows, creating activities without a
source requirement, and mixing business analysis with implementation details. This chapter avoids them by writing
the use case specifications first and drawing the activity and sequence diagrams from the specifications.

%% =========================================================================
\section{Use Case Specifications and Behavioural Models}\label{sec:2.2}

Following Assignment~03, I specify four customer use cases and two operational use cases in detail. The
four customer use cases are the two main ways of searching and the purchase flow: UC02 Search by Keyword (with its
voice variant UC03), UC04 Search by Image, UC07 Checkout and UC11 Cancel Order. The two operational use cases are
UC12 Manage Products and UC13 Update Order Status of the administrator. The use cases in which the AI advises the
customer directly (UC10, UC14) are specified in Chapter~3. Every specification uses the template of the
Lecture~03 slides: primary actor, supporting actor, goal, trigger, preconditions, postconditions, main flow and at
least three alternative or exception flows. The business rule ids (BR) are explained in Table~\ref{tab:rules}.

Each specification is followed directly by the diagram that shows its behaviour. An activity diagram shows the flow of work: initial node, activities, decisions, alternative branches, merges,
interactions with external parties, business rules and final node. Lecture~03 warns not to model every click but
the activities that mean something to the business. Họa. has six
activities: multimodal search (ACT-01), search by image (ACT-04), checkout (ACT-07), cancel order (ACT-11), manage
products (ACT-12) and update order status (ACT-13): every specified use case has its activity. Voice search is also shown as a sequence diagram, and cancellation depends on the state machine
of the order.

Three use cases also get an \emph{analysis sequence diagram} in the boundary--control--entity style of Lecture~03:
a \emph{boundary} object is the part of the system the actor talks to (a page), a \emph{control} object carries the
logic of one use case, and \emph{entity} objects are the analysis classes that hold the data (Section~\ref{sec:2.7}).
These diagrams say which object is responsible for each step, but nothing yet about classes of the code, method
signatures, storage or technology; the same scenarios return in Chapter~4 as design sequences with the real classes.

\newcommand{\specrow}[2]{\small\textbf{#1} & \small #2\\}

\needspace{10\baselineskip}
\subsection{UC02 Search by Keyword and UC03 Search by Voice}\label{sec:2.2.1}

\begin{longtable}{@{}P{3.0cm}P{12.0cm}@{}}
\caption{Specification of UC02 Search by Keyword (with the variant UC03 Search by Voice)}\label{tab:spec02}\\
\toprule
\textbf{Element} & \textbf{Description}\\
\midrule
\endfirsthead
\toprule
\textbf{Element} & \textbf{Description}\\
\midrule
\endhead
\specrow{Primary actor}{Customer (no sign-in needed)}
\specrow{Supporting actor}{Gemma 4 Service (understands the sentence)}
\specrow{Goal}{Find the paintings that match a free description, for example ``blue landscape paintings'' or ``Monet paintings under 100 dollars''.}
\specrow{Trigger}{The customer sends a search sentence (UC03: the customer speaks it).}
\specrow{Preconditions}{The catalogue has paintings for sale.}
\specrow{Postconditions}{The customer sees the paintings ordered by relevance, with the filters the system understood and the reason of each suggestion; no data is changed.}
\specrow{Main flow}{1. The customer enters a search sentence and sends it. 2. The system normalises the sentence (lower case, extra spaces removed). 3. The system asks the Gemma 4 Service to extract the conditions: style, subject, colour, artist, maximum price. 4. The system takes the candidate paintings that satisfy the hard conditions (artist, maximum price) and share keywords with the query. 5. The system scores and ranks the candidates. 6. The system shows the results, the conditions it understood and the reasons. 7. The customer opens a painting (UC05) or changes the sentence.}
\specrow{Voice variant (UC03)}{1a. The customer presses the microphone button and speaks. 1b. The system transcribes the audio with the Whisper model running on the same machine. 1c. The transcript appears in the search box; the customer corrects it if needed and confirms. Continue at step 2.}
\specrow{A1 -- AI does not answer}{At step 3 Gemma fails or times out: the system extracts the conditions with keyword rules (for example ``under 100 dollars'' $\rightarrow$ maximum price 100), says that it works in fallback mode and continues at step 4 (NFR02).}
\specrow{A2 -- No result}{At step 5 no candidate is left: the system shows an empty state and suggests removing a condition. No result is a valid result, not an error.}
\specrow{A3 -- Empty sentence}{At step 2 the sentence has only spaces or punctuation: the system asks for a search sentence and does not search.}
\specrow{A4 -- Speech cannot be transcribed}{At step 1b the Whisper server is down or the audio cannot be read: the system invites the customer to type the query (UC02) instead of showing an error.}
\specrow{Rules}{BR02 (sold-out prints can be viewed but not bought), BR11 (the AI never changes price or stock).}
\specrow{Requirements}{FR01, FR02, NFR02, NFR07, NFR08}
\bottomrule
\end{longtable}

The behaviour of UC02 and UC03 is the activity of multimodal search (ACT-01), which also covers the image
branch of UC04.

\begin{figure}[H]
\centering
\vpfig{c2_activity_search}{0.62\textheight}
\caption[Activity of multimodal search (ACT-01)]{Activity of multimodal search (ACT-01): three kinds of input meet in one common query, with a fallback branch when the AI service fails. Visual Paradigm diagram \vpsrc{c2_activity_search}.}\label{fig:c2_activity_search}\label{fig:actsearch}
\end{figure}

Figure~\ref{fig:actsearch} shows the central principle of Assignment~06: \emph{different inputs $\rightarrow$
common query representation $\rightarrow$ candidate retrieval $\rightarrow$ ranking $\rightarrow$ results}. The
first decision depends on the kind of input; speech is transcribed into text and then follows the same path as
typed text, while an image is described as attributes. The decision ``AI answered?'' comes from NFR02: when the
Gemma service does not answer, the flow does not stop but switches to keyword rules and joins again at candidate
retrieval.

Figure~\ref{fig:c2_bce_search} is the analysis sequence of UC02. The customer only talks to the search page
(\emph{boundary}); the search control (\emph{control}) interprets the request, asks the \code{Painting} entities for
candidates -- each painting knows its \code{Artist} -- and ranks them. The alternative fragment shows the two outcomes
of the specification: ranked results with their reasons, or an empty state with a hint (A2).

\begin{figure}[H]
\centering
\vpfig{c2_bce_search}{0.42\textheight}
\caption[Analysis sequence of search by keyword (boundary, control, entity)]{Analysis sequence of search by keyword (UC02): search page, search control and the Painting and Artist entities. Visual Paradigm diagram \vpsrc{c2_bce_search}.}\label{fig:c2_bce_search}
\end{figure}

\needspace{10\baselineskip}
\subsection{Interaction Model: Search by Voice}\label{sec:2.8}

Assignment~06~\cite{slidesA06} asks for a sequence diagram of the scenario ``the customer searches products by voice'' with exactly
seven participants: Customer, SearchUI, SpeechService, QueryService, SearchService, ProductRepository and
RankingService. Figure~\ref{fig:seqvoice} keeps the seven participants and the eight steps of the assignment:
(1) the customer speaks; (2) SearchUI sends the audio to SpeechService; (3) SpeechService returns the text;
(4) QueryService builds the common query; (5) SearchService retrieves candidates; (6) ProductRepository accesses the
painting data; (7) RankingService ranks; (8) SearchUI displays the results.

\begin{figure}[H]
\centering
\vpfig{c2_seq_voice}{0.55\textheight}
\caption[Sequence of search by voice (Assignment 06)]{Sequence of search by voice with exactly the seven participants of Assignment~06. Visual Paradigm diagram \vpsrc{c2_seq_voice}.}\label{fig:c2_seq_voice}\label{fig:seqvoice}
\end{figure}

Họa. differs from the simulation in the assignment only at steps 2--3: instead of returning a fixed string,
SpeechService calls a Whisper speech recognition model running on the same machine, so the audio never leaves the
machine (NFR08). The reasons for this choice and how it is measured are given in Chapter~3. Every step
after step 3 is unchanged: the transcript takes exactly the path of a typed sentence, as FR02 requires.

\needspace{10\baselineskip}
\subsection{UC04 Search by Image}\label{sec:2.2.2}

\begin{longtable}{@{}P{3.0cm}P{12.0cm}@{}}
\caption{Specification of UC04 Search by Image}\label{tab:spec04}\\
\toprule
\textbf{Element} & \textbf{Description}\\
\midrule
\endfirsthead
\toprule
\textbf{Element} & \textbf{Description}\\
\midrule
\endhead
\specrow{Primary actor}{Customer}
\specrow{Supporting actor}{Gemma 4 Service (describes the image)}
\specrow{Goal}{Find the prints whose style, subject and colour are close to a photo the customer has, for example a photo of a painting seen in a caf\'e.}
\specrow{Trigger}{The customer chooses an image file and starts the search (``show me paintings similar to this image'').}
\specrow{Preconditions}{The file is JPEG, PNG or WebP and not larger than 5 MB.}
\specrow{Postconditions}{The customer sees the closest paintings, each with a similarity score and a reason; the query image is not stored.}
\specrow{Main flow}{1. The customer chooses an image. 2. The page shows a preview and checks the format and size. 3. The system sends the image to Gemma 4 to describe its style, subject and dominant colour. 4. The system turns the description into a feature vector. 5. The system compares the vector of the image with the vector of every painting by cosine similarity. 6. The system ranks and shows the closest paintings. 7. The customer opens a painting.}
\specrow{A1 -- Wrong format or too large}{At step 2: the system refuses at once, states the allowed formats and size and does not call Gemma (NFR03).}
\specrow{A2 -- AI unavailable}{At step 3 Gemma 4 fails or does not answer in time: the system describes the photo locally by its dominant colour and continues with that description; the results are weaker, and every other function keeps working (NFR02).}
\specrow{A3 -- Not a painting}{At step 3 Gemma can describe only part of the image (for example a photo of an object): the system uses the attributes it has, the similarity is lower and the results are labelled as suggestions.}
\specrow{A4 -- Nothing close}{At step 6 all scores are low: the system shows an empty state instead of a noisy list.}
\specrow{Rules}{BR10 (image format and size), BR11.}
\specrow{Requirements}{FR03, NFR02, NFR03}
\bottomrule
\end{longtable}

The behaviour of UC04 is the image branch of the activity in Figure~\ref{fig:actsearch}: the photo is checked
(BR10), described as attributes and then follows the common path of every query. Its sequence with Gemma~4 is
analysed in Chapter~3 (Figure~\ref{fig:c3_seq_image}).

Figure~\ref{fig:c2_activity_image} draws UC04 on its own (ACT-04). The first decision applies BR10 and NFR03: a file
that is not JPEG, PNG or WebP, or is larger than 5~MB, is refused at once and Gemma~4 is never called. For a valid
photo Gemma~4 describes the style, subject and dominant colour with the shop's vocabulary; the second decision
applies NFR02: on an error or a timeout the system reads the average colour of the photo locally and names the
closest catalogue colour. Both paths then build the same query vector, retrieve the paintings by cosine similarity and
rank them; if nothing is close the customer sees an empty state (A4).

\begin{figure}[H]
\centering
\vpfig{c2_activity_image}{0.62\textheight}
\caption[Activity of search by image (ACT-04)]{Activity of search by image (ACT-04) with swimlanes for the customer, the Họa. system and Gemma 4 Service: the file check (BR10) and the local colour fallback (NFR02). Visual Paradigm diagram \vpsrc{c2_activity_image}.}\label{fig:c2_activity_image}
\end{figure}

\needspace{10\baselineskip}
\subsection{UC07 Checkout}\label{sec:2.2.3}

\begin{longtable}{@{}P{3.0cm}P{12.0cm}@{}}
\caption{Specification of UC07 Checkout}\label{tab:spec07}\\
\toprule
\textbf{Element} & \textbf{Description}\\
\midrule
\endfirsthead
\toprule
\textbf{Element} & \textbf{Description}\\
\midrule
\endhead
\specrow{Primary actor}{Customer (signed in)}
\specrow{Supporting actor}{Payment Service (confirms card payments)}
\specrow{Goal}{Create a valid order from the cart.}
\specrow{Trigger}{The customer confirms the order.}
\specrow{Preconditions}{The customer is signed in; the cart has at least one line.}
\specrow{Postconditions}{An order in status PLACED exists with all its lines at frozen prices; stock has been decreased; the cart is empty; one payment attempt is recorded.}
\specrow{Main flow}{1. The customer opens the cart. 2. The system shows the prints, quantities and the provisional total. 3. The customer confirms, choosing a delivery address (a saved one, or a new one typed in: A6) and a payment method. 4. The system reads every product again and checks that enough stock is left. 5. The system creates the order lines at the current catalogue price (frozen price). 6. The system computes the total: subtotal minus the VIP discount plus the shipping fee. 7. The system records the payment: cash is collected on delivery; a card needs an authorisation code from the payment service. 8. The system saves the order, decreases the stock and empties the cart. 9. The system shows the order id and the status PLACED.}
\specrow{A1 -- Empty cart}{At step 3: checkout is refused (BR03).}
\specrow{A2 -- Out of stock}{At step 4 one print is short: the system names it, asks the customer to change the cart and creates no order (BR02).}
\specrow{A3 -- Invalid address}{At step 3 the address does not belong to the customer: the system asks for one of the customer's own addresses.}
\specrow{A4 -- Card refused}{At step 7 there is no authorisation code: the payment is recorded as FAILED and the order is not considered paid (BR09).}
\specrow{A5 -- Price changed}{At step 5 the catalogue price differs from the price when the print was added: the order uses the current price and the customer sees the new total (BR04).}
\specrow{A6 -- New address}{At step 3 none of the saved addresses fits: the customer types the house number, street, district and city. The system removes extra blanks, refuses an empty or too long part and saves the address for this customer; the new address is selected and the flow goes on at step 3.}
\specrow{Rules}{BR01, BR02, BR03, BR04, BR05, BR09}
\specrow{Requirements}{FR07, NFR05}
\bottomrule
\end{longtable}

\begin{figure}[H]
\centering
\vpfig{c2_activity_checkout}{0.66\textheight}
\caption[Activity of checkout with swimlanes (ACT-07)]{Activity of checkout with swimlanes (ACT-07): customer, Họa. system and database. Visual Paradigm diagram \vpsrc{c2_activity_checkout}.}\label{fig:c2_activity_checkout}\label{fig:actcheckout}
\end{figure}

Figure~\ref{fig:actcheckout} divides the work into swimlanes: the customer views the cart and confirms; the system
checks, freezes prices, computes the total and records the payment; the database saves the order, decreases the
stock and empties the cart. Each decision on the diagram corresponds to a business rule (Table~\ref{tab:decisions}).
The order of the steps matters: the stock check and the price freeze come before saving, and the order is saved
before the stock is decreased and the cart is emptied, so that a state ``stock decreased but no order'' can never
exist.

Figure~\ref{fig:c2_bce_checkout} gives the same scenario as an analysis sequence. The checkout page is the boundary
and the checkout control carries the use case. A loop checks the stock and current price of each cart line in
\code{Painting}; the alternative then either refuses and names the missing item (BR02) or lets \code{Order} create
its \code{OrderItem}s at the frozen price and compute its own total, before the payment is recorded, the stock is
reduced and the cart is emptied. Chapter~4 shows the same flow through \code{CheckoutController},
\code{CheckoutService} and the repositories (Figure~\ref{fig:c44_checkout_seq}).

\begin{figure}[H]
\centering
\vpfig{c2_bce_checkout}{0.50\textheight}
\caption[Analysis sequence of checkout (boundary, control, entity)]{Analysis sequence of checkout (UC07): checkout page, checkout control and the Cart, Painting, Order, OrderItem and Payment entities. Visual Paradigm diagram \vpsrc{c2_bce_checkout}.}\label{fig:c2_bce_checkout}
\end{figure}

\needspace{10\baselineskip}
\subsection{UC11 Cancel Order}\label{sec:2.2.4}

\begin{longtable}{@{}P{3.0cm}P{12.0cm}@{}}
\caption{Specification of UC11 Cancel Order}\label{tab:spec11}\\
\toprule
\textbf{Element} & \textbf{Description}\\
\midrule
\endfirsthead
\toprule
\textbf{Element} & \textbf{Description}\\
\midrule
\endhead
\specrow{Primary actor}{Customer (signed in)}
\specrow{Supporting actor}{None}
\specrow{Goal}{Cancel one of the customer's own orders while it is not shipped.}
\specrow{Trigger}{The customer chooses ``Cancel order'' on an order (extends UC08 View Order).}
\specrow{Preconditions}{The customer is signed in.}
\specrow{Postconditions}{The order is CANCELLED and cannot move any more.}
\specrow{Main flow}{1. The customer opens the list of orders. 2. The system shows the customer's orders and their statuses. 3. The customer chooses an order and asks to cancel it. 4. The system checks that the order belongs to the signed-in customer. 5. The system checks that the status allows cancellation (PLACED or PROCESSING). 6. The system sets the order to CANCELLED and shows the new status.}
\specrow{A1 -- Already shipped}{At step 5 the order is SHIPPED or DELIVERED (for example order 20261001): the system refuses and explains that only orders not yet shipped can be cancelled (BR06).}
\specrow{A2 -- Another customer's order}{At step 4: the system answers as if the order did not exist and reveals nothing (BR08).}
\specrow{A3 -- Already cancelled}{At step 5 the order is already CANCELLED: the system says so and changes nothing.}
\specrow{Open question}{Restoring stock and refunding on cancellation are recorded for the owner to decide; the demo only changes the order status.}
\specrow{Rules}{BR06, BR08}
\specrow{Requirements}{FR11, NFR01}
\bottomrule
\end{longtable}

Figure~\ref{fig:c2_activity_cancel} shows the cancellation (ACT-11). The customer's identity comes from the
signed-in session, never from the request. The first decision applies BR08 and NFR01: an order that does not exist and
an order of another customer receive the same answer, ``not found in your account'', so nothing leaks. The second
decision applies BR06: only a PLACED or PROCESSING order can be cancelled; a SHIPPED, DELIVERED or already CANCELLED
order is refused and nothing changes (A1, A3). Otherwise the order becomes CANCELLED, is saved and the customer sees
the new status (FR11).

\begin{figure}[H]
\centering
\vpfig{c2_activity_cancel}{0.45\textheight}
\caption[Activity of cancel order (ACT-11)]{Activity of order cancellation (ACT-11) with swimlanes for the customer, the Họa. system and the database: the ownership check (BR08) and the status check (BR06). Visual Paradigm diagram \vpsrc{c2_activity_cancel}.}\label{fig:c2_activity_cancel}
\end{figure}

Figure~\ref{fig:c2_bce_cancel} is the analysis sequence of the same use case: the order page receives the request,
the cancel control finds the \code{Order} and checks its owner and status, and the alternative fragment covers the
four outcomes of the specification -- another customer's order (answered as not found, BR08), a shipped order
(refused, BR06), an order already cancelled (A3) and a PLACED or PROCESSING order, which is told to \code{cancel()}.

\begin{figure}[H]
\centering
\vpfig{c2_bce_cancel}{0.42\textheight}
\caption[Analysis sequence of cancel order (boundary, control, entity)]{Analysis sequence of order cancellation (UC11): order page, cancel control and the Order entity. Visual Paradigm diagram \vpsrc{c2_bce_cancel}.}\label{fig:c2_bce_cancel}
\end{figure}

Whether an order may still be cancelled depends on its state, which is modelled next.

\needspace{10\baselineskip}
\subsection{State Model of the Order}\label{sec:2.9}

An order has a clear life cycle and many rules depend on its state, so it is modelled with a UML state machine~\cite{uml}
(Figure~\ref{fig:stateorder}). Payment has its own life cycle (PENDING $\rightarrow$ SUCCEEDED or FAILED), which
is not the life cycle of the order: a failed payment does not cancel the order by itself, and an order being
shipped does not mean that the money has been collected.

\begin{figure}[H]
\centering
\vpfig{c2_state_order}{0.45\textheight}
\caption[State machine of the order]{State machine of \texttt{Order}: the values of \texttt{OrderStatus} and the valid transitions. Visual Paradigm diagram \vpsrc{c2_state_order}.}\label{fig:c2_state_order}\label{fig:stateorder}
\end{figure}

\begin{longtable}{@{}P{3.6cm}P{3.6cm}P{3.6cm}P{4.2cm}@{}}
\caption{State transitions of the order}\label{tab:transitions}\\
\toprule
\textbf{Transition} & \textbf{Event} & \textbf{Actor} & \textbf{Condition (rule)}\\
\midrule
\endfirsthead
\toprule
\textbf{Transition} & \textbf{Event} & \textbf{Actor} & \textbf{Condition (rule)}\\
\midrule
\endhead
\small (initial) $\rightarrow$ PLACED & \small Checkout succeeded & \small Customer & \small BR02, BR03\\
\small PLACED $\rightarrow$ PROCESSING & \small Warehouse accepts the order & \small Administrator & \small BR07\\
\small PROCESSING $\rightarrow$ SHIPPED & \small Handed to the carrier & \small Administrator & \small BR07\\
\small SHIPPED $\rightarrow$ DELIVERED & \small Delivery confirmed & \small Administrator & \small BR07\\
\small PLACED $\rightarrow$ CANCELLED & \small Customer cancels & \small Customer & \small BR06, BR08\\
\small PROCESSING $\rightarrow$ CANCELLED & \small Customer cancels & \small Customer & \small BR06, BR08\\
\bottomrule
\end{longtable}

DELIVERED and CANCELLED are final states, with no backward transition and no jump; this is what the consistency
check of Table~\ref{tab:consistency} compares with the operations of class \code{Order}.

\needspace{10\baselineskip}
\subsection{UC12 Manage Products and UC13 Update Order Status}\label{sec:2.2.5}

\begin{longtable}{@{}P{3.0cm}P{12.0cm}@{}}
\caption{Specification of UC12 Manage Products (operational)}\label{tab:spec12}\\
\toprule
\textbf{Element} & \textbf{Description}\\
\midrule
\endfirsthead
\toprule
\textbf{Element} & \textbf{Description}\\
\midrule
\endhead
\specrow{Primary actor}{Administrator}
\specrow{Supporting actor}{None}
\specrow{Goal}{Keep the catalogue right: add new paintings, change price and stock.}
\specrow{Trigger}{The administrator opens the product management page and edits it.}
\specrow{Preconditions}{Signed in with an administrator session.}
\specrow{Postconditions}{The catalogue is updated; a new price applies to later orders, old orders keep their frozen prices.}
\specrow{Main flow}{1. The administrator opens the list of paintings. 2. The system shows the paintings with price and stock. 3. The administrator changes the price or stock of a painting, or enters a new painting with its category (artist, style, subject, colour, size, image). 4. The system validates the data and checks that the category accepts a painting. 5. The system saves the change. 6. The system shows the updated list.}
\specrow{A1 -- Invalid data}{At step 4 a negative price or stock, or a missing mandatory field: the system refuses and names the wrong field.}
\specrow{A2 -- No permission}{At step 1 the session is not an administrator session: access is refused.}
\specrow{A3 -- Print in a cart}{The price changes while the print is in a customer's cart: the cart keeps the price of the moment it was added, and the order will use the new price at checkout (BR04).}
\specrow{A4 -- Wrong category}{At step 4 the chosen category holds another kind of product (a painting sent to Frames): the system refuses and writes nothing.}
\specrow{Requirements}{FR12, NFR01}
\bottomrule
\end{longtable}

\begin{longtable}{@{}P{3.0cm}P{12.0cm}@{}}
\caption{Specification of UC13 Update Order Status (operational)}\label{tab:spec13}\\
\toprule
\textbf{Element} & \textbf{Description}\\
\midrule
\endfirsthead
\toprule
\textbf{Element} & \textbf{Description}\\
\midrule
\endhead
\specrow{Primary actor}{Administrator}
\specrow{Supporting actor}{None}
\specrow{Goal}{Move an order through processing and delivery.}
\specrow{Trigger}{The warehouse completes a step (order accepted, handed to the carrier, delivery confirmed).}
\specrow{Preconditions}{Signed in as administrator; the order exists.}
\specrow{Postconditions}{The order has its new status; the customer sees it when viewing or looking up the order.}
\specrow{Main flow}{1. The administrator opens the list of all orders. 2. The system shows the orders and their statuses. 3. The administrator chooses the next status of an order. 4. The system checks that the transition is valid. 5. The system saves and shows the new status.}
\specrow{A1 -- Skipped step}{At step 4 a jump PLACED $\rightarrow$ DELIVERED: refused because the order of statuses is wrong (BR07).}
\specrow{A2 -- Cancelled order}{The order is CANCELLED: no transition is possible.}
\specrow{A3 -- No permission}{The session is not an administrator session: refused.}
\specrow{Requirements}{FR13, NFR01}
\bottomrule
\end{longtable}

UC13 moves an order only along the transitions of Figure~\ref{fig:stateorder}.

Figure~\ref{fig:c2_activity_products} shows UC12 (ACT-12). The back office opens only for an administrator session
(A2). The administrator either changes the price or stock of a product or enters a new painting with its category;
both meet the same validation, which refuses a negative value or a missing field and names it (A1). A new painting
must also pass \code{Category.accepts}: a category that holds another kind of product is refused and nothing is
written (A4). An accepted painting is saved as its product row, its painting row and its feature vector in one
transaction.

\begin{figure}[H]
\centering
\vpfig{c2_activity_products}{0.66\textheight}
\caption[Activity of manage products (ACT-12)]{Activity of product management (ACT-12) with swimlanes for the administrator, the Họa. system and the database: the session check, the validation, the category check of a new painting and its save in one transaction. Visual Paradigm diagram \vpsrc{c2_activity_products}.}\label{fig:c2_activity_products}
\end{figure}

Figure~\ref{fig:c2_activity_admin} shows UC13 (ACT-13). The administrator selects an order and chooses its next
status; the system checks the request against the state machine of Figure~\ref{fig:stateorder}: only the next step of
PLACED, PROCESSING, SHIPPED, DELIVERED is allowed, and CANCELLED and DELIVERED are final. A skipped step, a step back
or a move from a final status is refused (BR07); a valid transition is saved and shown, and the customer sees it the
next time the order is opened.

\begin{figure}[H]
\centering
\vpfig{c2_activity_admin}{0.45\textheight}
\caption[Activity of update order status (ACT-13)]{Activity of order status update (ACT-13) with swimlanes for the administrator, the Họa. system and the database: the transition check of the order state machine (BR07). Visual Paradigm diagram \vpsrc{c2_activity_admin}.}\label{fig:c2_activity_admin}
\end{figure}

\subsection{Summary of alternative flows}

Lecture~03 insists that the normal flow is not enough. Table~\ref{tab:altsummary} collects the most important
alternative situations of the specifications above as \emph{situation $\rightarrow$ system behaviour}, like the
Place Order table of the Lecture~03 slides.

\begin{longtable}{@{}P{4.6cm}P{7.8cm}P{2.6cm}@{}}
\caption{Alternative situations and system behaviour}\label{tab:altsummary}\\
\toprule
\textbf{Situation} & \textbf{System behaviour} & \textbf{Use case}\\
\midrule
\endfirsthead
\toprule
\textbf{Situation} & \textbf{System behaviour} & \textbf{Use case}\\
\midrule
\endhead
\small A print is out of stock & \small Name the print, ask to change the cart, create no order & \small UC07\\
\small The card is refused & \small Record a failed payment; the order is not paid & \small UC07\\
\small Cancelling a shipped order & \small Refuse and explain the cancellation rule & \small UC11\\
\small Asking for another customer's order & \small Answer ``not found'', reveal nothing & \small UC09, UC11\\
\small AI fails during text search & \small Use keyword rules and say so & \small UC02, UC03\\
\small Whisper server is down & \small Invite the customer to type instead & \small UC03\\
\small Image of wrong format or too large & \small Refuse before calling the AI & \small UC04, UC14\\
\small Status moved in the wrong order & \small Refuse and keep the old status & \small UC13\\
\bottomrule
\end{longtable}

%% =========================================================================
\section{Data Flow Diagram (Level 1)}\label{sec:dfd1}

The use cases say what the actors want; the data flow diagram says how the information moves through the system and
where it is kept. Level~1 opens process~0 of the context diagram (Figure~\ref{fig:context}) into four processes and
three data stores (Figure~\ref{fig:c2_dfd1}):
\begin{itemize}
\item \textbf{1 Find paintings and answer questions} (UC01--UC05, UC10, UC14) receives the customer's query or question,
asks Gemma~4 to interpret it, reads the paintings (D1) and, for the assistant, the customer's own orders (D3).
\item \textbf{2 Shop and check out} (UC06, UC07) keeps the cart (D2), reads price and stock and decreases the stock (D1),
asks the payment service for an authorisation and writes the new order with its payment (D3).
\item \textbf{3 Follow and cancel orders} (UC08, UC09, UC11) reads the orders and writes the status CANCELLED (D3).
\item \textbf{4 Administer catalogue and orders} (UC12, UC13) writes painting changes (D1) and new order statuses (D3).
\end{itemize}
The diagram is \emph{balanced} with level~0: every flow of the context diagram appears again between the same
external entity and one of the processes, and no new external entity appears. Every data store has a flow in and a
flow out. To keep the lines from crossing, the customer and D1 are drawn twice, as DFDs allow.

\landscapefig{c2_dfd1}{Data flow diagram, level 1}{Data flow diagram of Họa., level 1: four processes, three data stores and the external entities of the context diagram. \trialshot{data flow diagrams}. Diagram \vpsrc{c2_dfd1}.}{\label{fig:c2_dfd1}}

%% =========================================================================
\section{Business Rules}\label{sec:2.3}

A business rule is a condition set by the organisation, independent of the user interface. In an activity model
every decision point must come from a business rule (Lecture~03: ``business rules drive decisions'').
Table~\ref{tab:rules} lists the rules of Họa.; the BR ids are used in the use case specifications, the activity
diagrams and the state machine.

\begin{longtable}{@{}P{1.3cm}P{10.0cm}P{3.7cm}@{}}
\caption{Business rules of Họa.}\label{tab:rules}\\
\toprule
\textbf{ID} & \textbf{Rule} & \textbf{Decision point}\\
\midrule
\endfirsthead
\toprule
\textbf{ID} & \textbf{Rule} & \textbf{Decision point}\\
\midrule
\endhead
\small BR01 & \small A customer has at most one active cart. & \small Add to cart\\
\small BR02 & \small A print cannot be ordered when its stock is lower than the quantity; a sold-out print can still be viewed. & \small ``In stock?''\\
\small BR03 & \small A valid order has at least one line; an empty cart cannot be checked out. & \small ``Cart has items?''\\
\small BR04 & \small The price on an order is the catalogue price at checkout time and does not change when the catalogue price changes later. & \small Create order lines\\
\small BR05 & \small A VIP customer gets the discount rate of the tier; the total is never negative. & \small Compute total\\
\small BR06 & \small An order can be cancelled only while it is PLACED or PROCESSING. & \small ``Cancellation allowed?''\\
\small BR07 & \small The order status only moves PLACED $\rightarrow$ PROCESSING $\rightarrow$ SHIPPED $\rightarrow$ DELIVERED. & \small ``Valid transition?''\\
\small BR08 & \small Customers view, look up and cancel only their own orders; the identity comes from the signed-in session. & \small ``Customer's order?''\\
\small BR09 & \small A card payment succeeds only with an authorisation code; the full card number is never stored. & \small ``Payment succeeded?''\\
\small BR10 & \small Uploaded images are JPEG, PNG or WebP and at most 5 MB. & \small ``Valid image?''\\
\small BR11 & \small AI suggestions never change price, stock or access rights. & \small Every use case with AI\\
\bottomrule
\end{longtable}

Each decision in the activity diagrams of Section~\ref{sec:2.2} comes from one of these rules
(Table~\ref{tab:decisions}).

\begin{longtable}{@{}P{4.2cm}P{8.3cm}P{2.5cm}@{}}
\caption{Decision points of the activities and their business rules}\label{tab:decisions}\\
\toprule
\textbf{Decision} & \textbf{Business rule} & \textbf{Activity}\\
\midrule
\endfirsthead
\toprule
\textbf{Decision} & \textbf{Business rule} & \textbf{Activity}\\
\midrule
\endhead
\small Kind of input? & \small Text, voice and image all become one common query (Assignment~06) & \small ACT-01\\
\small AI answered? & \small NFR02: if not, use keyword rules & \small ACT-01\\
\small Valid image? & \small BR10 & \small ACT-01\\
\small Cart has items? & \small BR03 & \small ACT-07\\
\small In stock? & \small BR02 & \small ACT-07\\
\small Payment succeeded? & \small BR09 & \small ACT-07\\
\small Image accepted? & \small BR10, NFR03: JPEG, PNG or WebP, at most 5~MB & \small ACT-04\\
\small Description received? & \small NFR02: if not, describe the photo locally by its colour & \small ACT-04\\
\small Customer's order? & \small BR08: another customer's order is answered as not found & \small ACT-11\\
\small Cancellation allowed? & \small BR06: only PLACED or PROCESSING & \small ACT-11\\
\small Valid transition? & \small BR07: only the next status of the state machine & \small ACT-13\\
\small Valid product data? & \small No negative price or stock, no missing field (UC12 A1) & \small ACT-12\\
\small Category accepts the painting? & \small The detail kind of the category must be PAINTING (UC12 A4, NFR11) & \small ACT-12\\
\bottomrule
\end{longtable}

%% =========================================================================
\section{Domain Model}\label{sec:2.7}

Lecture~04 moves from behaviour to structure~\cite{slidesL04,larman,evans}: \emph{requirements + behaviour $\Rightarrow$ domain concepts}. The
concepts appear naturally when the scenarios are read again (Table~\ref{tab:behavior2concept}).

\begin{longtable}{@{}P{6.4cm}P{8.6cm}@{}}
\caption{From business behaviour to domain concepts}\label{tab:behavior2concept}\\
\toprule
\textbf{Business behaviour} & \textbf{Domain concepts}\\
\midrule
\endfirsthead
\toprule
\textbf{Business behaviour} & \textbf{Domain concepts}\\
\midrule
\endhead
\small A customer adds a print to the cart & \small \code{Customer}, \code{Cart}, \code{CartItem}, \code{Product}\\
\small The customer changes a quantity & \small \code{CartItem.quantity}\\
\small The customer places an order & \small \code{Cart}, \code{Order}\\
\small An order contains the purchased prints & \small \code{Order}, \code{OrderItem}, \code{Product}\\
\small The purchase price must be kept & \small \code{OrderItem.unitPrice}\\
\small A print reproduces a work of one artist & \small \code{Painting}, \code{Artist}\\
\small The customer buys a frame with the print & \small \code{Frame}\\
\small The shop sells two lines, prints and frames, and the owner manages each line apart & \small \code{Category}\\
\small The order is delivered to an address & \small \code{Shipping}, \code{Address}\\
\small The customer pays cash or by card & \small \code{Payment}, \code{PayCash}, \code{PayCredit}\\
\small A VIP customer gets a discount & \small \code{CustomerVIP}\\
\bottomrule
\end{longtable}

\begin{figure}[H]
\centering
\vpfig{c2_domain}{0.6\textheight}
\caption[Domain model of Họa.]{Domain model of Họa. drawn as the analysis class diagram: \nAnalysisClasses{} classes with their attributes and the kind of each relationship; names and multiplicities are added in the data model (Section~\ref{sec:2.dm.2}). Visual Paradigm diagram \vpsrc{c2_domain}.}\label{fig:c2_domain}\label{fig:domain}
\end{figure}

Figure~\ref{fig:domain} is the analysis class diagram: each class shows its attributes, and each line shows only the
kind of relationship -- association, composition (filled diamond), aggregation (hollow diamond), generalization
(hollow triangle) and the «create» dependency of checkout. The relationships get their names and multiplicities one
step later, in the logical data model (Section~\ref{sec:2.dm.2}), and their direction and methods in the design
(Chapter~4). The diagram contains business concepts only. Following the Lecture~04 slides, the analysis model does
\emph{not} contain CheckoutController, CheckoutService, CartRepository, OrderRepository, REST APIs, database tables
or layers; they answer the question ``how'' and belong to Chapter~4. The conceptual responsibilities
of the main classes are listed in Table~\ref{tab:responsibility}.

\begin{longtable}{@{}P{2.8cm}P{12.2cm}@{}}
\caption{Conceptual responsibilities in the analysis model}\label{tab:responsibility}\\
\toprule
\textbf{Concept} & \textbf{Responsibility}\\
\midrule
\endfirsthead
\toprule
\textbf{Concept} & \textbf{Responsibility}\\
\midrule
\endhead
\small \code{Customer} & \small Owns a cart, places orders, has one name and at least one address.\\
\small \code{Product} & \small Something the gallery sells: a print (\code{Painting}) or a frame (\code{Frame}); has a price and a stock.\\
\small \code{Painting} & \small A print of one painting by one artist; carries the attributes used for search: style, subject, colour, medium, year, size.\\
\small \code{Cart} & \small Holds the current shopping state.\\
\small \code{CartItem} & \small Records one product and the intended quantity in the cart, with the price when added.\\
\small \code{Order} & \small Represents a confirmed transaction with a status and a total.\\
\small \code{OrderItem} & \small Records the product, the purchased quantity and the transaction price.\\
\small \code{Shipping}, \code{Payment} & \small Deliver one order to one address; record the payment attempts of an order.\\
\bottomrule
\end{longtable}

\subsection{CartItem is not OrderItem}

The two classes have similar attributes but two different meanings (Lecture~04 slides): \code{CartItem} is
temporary shopping information, \code{OrderItem} is historical transaction information.

\begin{longtable}{@{}P{3.2cm}P{5.9cm}P{5.9cm}@{}}
\caption{CartItem compared with OrderItem}\label{tab:cartvsorder}\\
\toprule
\textbf{Criterion} & \textbf{\code{CartItem}} & \textbf{\code{OrderItem}}\\
\midrule
\endfirsthead
\toprule
\textbf{Criterion} & \textbf{\code{CartItem}} & \textbf{\code{OrderItem}}\\
\midrule
\endhead
\small Purpose & \small Shopping & \small Transaction\\
\small Lifetime & \small Temporary & \small Permanent\\
\small Belongs to & \small \code{Cart} & \small \code{Order}\\
\small Quantity & \small Intended quantity (\code{quantity}) & \small Purchased quantity (\code{quantity})\\
\small Price & \small Price when added to the cart (\code{cartPrice}) & \small Transaction price (\code{unitPrice})\\
\small After checkout & \small The cart is emptied & \small Stays in the order for ever\\
\bottomrule
\end{longtable}

\subsection{Checkout is a transformation}

The central transformation of the domain model is \emph{Cart} $\xrightarrow{\text{checkout}}$ \emph{Order}. For
example, the cart of customer C001 holds \{1 print of ``Water Lilies'' by Monet, 1 ``Natural oak frame 40 x 60''\};
after a successful checkout the new order holds \{OrderItem$_1$: ``Water Lilies'' $\times 1$, OrderItem$_2$: oak
frame $\times 1$\} at frozen prices, and \code{Cart.items} $= \varnothing$. For every line $i$:
\[
\mathrm{Product}(\mathrm{CartItem}_i) = \mathrm{Product}(\mathrm{OrderItem}_i), \qquad
\mathrm{Quantity}(\mathrm{CartItem}_i) = \mathrm{Quantity}(\mathrm{OrderItem}_i).
\]
Because this is a transformation and not a permanent link, the relationship between \code{Cart} and \code{Order}
is drawn as a «create» dependency, not as an association.


%% =========================================================================
\section{Data Model}\label{sec:2.dm}

The domain model says which concepts the shop works with; the data model says which of them must be kept and how.
Following Test~01, it is drawn at three levels, each adding detail to the one before (Table~\ref{tab:levels42}). The
conceptual and logical levels belong to the analysis and are drawn here; the physical level, with tables, data types
and keys, is a design decision and is drawn in Section~\ref{sec:4.2}, where the database is generated from it.

\begin{longtable}{@{}P{2.3cm}P{8.6cm}P{4.1cm}@{}}
\caption{The three levels of the data model}\label{tab:levels42}\\
\toprule
\textbf{Level} & \textbf{What it shows} & \textbf{Figure}\\
\midrule
\endfirsthead
\toprule
\textbf{Level} & \textbf{What it shows} & \textbf{Figure}\\
\midrule
\endhead
\small Conceptual & \small The entities by name and the named links between them; no attributes, no multiplicities & \small Figure~\ref{fig:c42_conceptual}\\
\small Logical & \small + attributes; every class «ORM-Persistable»; every link named, with multiplicities at both ends (1, 0..1, 0..*, 1..*); generalizations turned into \emph{is a} links & \small Figure~\ref{fig:c42_orm}\\
\small Physical & \small + tables and columns with SQL data types; primary, foreign and unique keys; each link becomes a foreign key with its name and multiplicities & \small Figure~\ref{fig:c42_tables} (Section~\ref{sec:4.2})\\
\bottomrule
\end{longtable}

\needspace{10\baselineskip}
\subsection{Conceptual Data Model}\label{sec:2.dm.1}

The conceptual data model (Figure~\ref{fig:c42_conceptual}) is the roughest picture: the \nAnalysisClasses{}
concepts of the domain model as entities, shown by name only, and the links between them with a name that reads as a
sentence (a customer \emph{owns} a cart, \emph{places} orders; an order \emph{contains} order items and \emph{is
paid by} payments; a category \emph{classifies} products). It differs from the domain model of Section~\ref{sec:2.7} in what it leaves out: no
multiplicities, no composition or aggregation, and no «create» dependency of checkout, because that dependency is
an operation, not something to store. The generalizations stay, because at this level a VIP customer is still
simply a kind of customer.

\landscapefig{c42_conceptual}{Conceptual data model of Họa.}{Conceptual data model: the entities by name and the named links between them. Visual Paradigm diagram \vpsrc{c42_conceptual}.}{\label{fig:c42_conceptual}}

\needspace{10\baselineskip}
\subsection{Logical Data Model}\label{sec:2.dm.2}

The logical data model (Figure~\ref{fig:c42_orm}) adds the attributes of every class (their meaning is
explained in Section~\ref{sec:4.1}) and follows the four steps of Test~01: every class is marked «ORM-Persistable», meaning that its objects are stored as records;
every association receives a name (its label) and multiplicities at both ends; every generalization becomes an
association \emph{is a} with multiplicity 1 at the parent and 0..1 at the child, because the relational model has
no inheritance; and the «create» dependency stays out. The multiplicities read in both directions: \emph{Customer
1 -- 0..* Order} says that a customer places zero or more orders and that every order belongs to exactly one
customer; \emph{Order 1 -- 1..* OrderItem} says that an order has at least one line; \emph{Customer 1 -- 0..1
CustomerVIP} says that a customer has at most one VIP record and every VIP record belongs to one customer.

\landscapefig{c42_orm}{Logical data model of Họa.}{Logical data model: attributes, classes marked «ORM-Persistable», named associations with multiplicities, generalizations turned into \emph{is a} associations 1 -- 0..1. Visual Paradigm diagram \vpsrc{c42_orm}.}{\label{fig:c42_orm}\label{fig:orm}}

\paragraph{Multiplicity.} Multiplicity must be modelled explicitly (Lecture~04 slides); in this essay it is written
at the logical level of the data model, together with the names of the relationships. Table~\ref{tab:mult} gives the
multiplicity of the main relationships and the business assumption behind each value.

\begin{longtable}{@{}P{3.7cm}P{2.4cm}P{8.9cm}@{}}
\caption{Multiplicities and business assumptions}\label{tab:mult}\\
\toprule
\textbf{Relationship} & \textbf{Multiplicity} & \textbf{Assumption}\\
\midrule
\endfirsthead
\toprule
\textbf{Relationship} & \textbf{Multiplicity} & \textbf{Assumption}\\
\midrule
\endhead
\small Customer -- Cart & \small 1 : 0..1 & \small At most one active cart per customer (BR01); a new customer may have none.\\
\small Cart -- CartItem & \small 1 : 0..* & \small An empty cart is valid.\\
\small CartItem -- Product & \small 0..* : 1 & \small Each cart line refers to exactly one product.\\
\small Order -- OrderItem & \small 1 : 1..* & \small An empty order makes no sense (BR03).\\
\small OrderItem -- Product & \small 0..* : 1 & \small Each order line refers to exactly one product.\\
\small Artist -- Painting & \small 1 : 0..* & \small Each print belongs to exactly one artist.\\
\small Order -- Payment & \small 1 : 0..* & \small An order can have several payment attempts.\\
\bottomrule
\end{longtable}

%% =========================================================================
\section{Requirement Derivation and Traceability}\label{sec:trace}\label{sec:1.12}

Lecture~02 stresses that every requirement needs a unique identifier and that the derivation chain \emph{problem
$\rightarrow$ need $\rightarrow$ functional requirement $\rightarrow$ actor $\rightarrow$ use case} must be kept
during the whole project. Table~\ref{tab:derivation} links the five problems to use cases; therefore the question
``why does this use case exist?'' is answered with a concrete business problem instead of a random list of
functions.

\begin{longtable}{@{}P{2.3cm}P{3.5cm}P{2.6cm}P{2.6cm}P{4.0cm}@{}}
\caption{Derivation chain problem -- need -- requirement -- actor -- use case}\label{tab:derivation}\\
\toprule
\textbf{Problem} & \textbf{Need} & \textbf{Requirements} & \textbf{Actors} & \textbf{Use cases}\\
\midrule
\endfirsthead
\toprule
\textbf{Problem} & \textbf{Need} & \textbf{Requirements} & \textbf{Actors} & \textbf{Use cases}\\
\midrule
\endhead
\small P1 hard to find paintings & \small Search by description, voice, photo, room & \small FR01--FR04, FR14, NFR08 & \small Customer, Gemma 4 Service & \small UC01--UC05, UC14\\
\small P2 sold-out prints sold & \small Accurate stock at checkout & \small FR05, FR07, FR12 & \small Customer, Administrator & \small UC05, UC07, UC12\\
\small P3 order status unknown & \small See and look up order status & \small FR08--FR11, FR13 & \small Customer, Administrator, Gemma 4 Service & \small UC08, UC09, UC10, UC11, UC13\\
\small P4 manual work & \small Administration tools & \small FR12, FR13 & \small Administrator & \small UC12, UC13\\
\small P5 risk of data leaks & \small Ownership checks, no card numbers, local voice & \small NFR01, NFR05, NFR08 & \small Customer, Payment Service & \small UC03, UC07, UC08, UC09, UC10\\
\bottomrule
\end{longtable}

Table~\ref{tab:traceuc} in Appendix~C is the two-way traceability matrix between requirements and use cases: every
requirement is realised by at least one use case and every use case has at least one source requirement. The trace
continues to scenarios and activities (Section~\ref{sec:2.5}), and to classes, tables and code in Section~4.6 and
Appendix~B.

%% =========================================================================
\section{Traceability: Requirement -- Use Case -- Scenario -- Activity}\label{sec:2.5}

Every important functional model must be traceable (Lecture~03): from a requirement one must find the use case,
the main scenario and the activity that show it.

\begin{longtable}{@{}P{1.5cm}P{3.6cm}P{5.6cm}P{4.3cm}@{}}
\caption{Trace FR $\rightarrow$ UC $\rightarrow$ scenario $\rightarrow$ activity}\label{tab:tracefr}\\
\toprule
\textbf{Req.} & \textbf{Use case} & \textbf{Main scenario} & \textbf{Behaviour model}\\
\midrule
\endfirsthead
\toprule
\textbf{Req.} & \textbf{Use case} & \textbf{Main scenario} & \textbf{Behaviour model}\\
\midrule
\endhead
\small FR01 & \small UC02 Search by Keyword & \small Main flow of Table~\ref{tab:spec02} & \small ACT-01 (Figure~\ref{fig:actsearch})\\
\small FR02 & \small UC03 Search by Voice & \small Voice variant of Table~\ref{tab:spec02} & \small ACT-01, sequence of Figure~\ref{fig:seqvoice}\\
\small FR03 & \small UC04 Search by Image & \small Main flow of Table~\ref{tab:spec04} & \small ACT-04 (Figure~\ref{fig:c2_activity_image}), sequence of Figure~\ref{fig:c3_seq_image}\\
\small FR07 & \small UC07 Checkout & \small Main flow of Table~\ref{tab:spec07} & \small ACT-07 (Figure~\ref{fig:actcheckout})\\
\small FR10 & \small UC10 Ask Shopping Assistant & \small Specification of UC10 (Chapter~3) & \small Sequence of Figure~\ref{fig:c3_seq_assistant}\\
\small FR11 & \small UC11 Cancel Order & \small Main flow of Table~\ref{tab:spec11} & \small ACT-11 (Figure~\ref{fig:c2_activity_cancel}), state machine\\
\small FR12 & \small UC12 Manage Products & \small Main flow of Table~\ref{tab:spec12} & \small ACT-12 (Figure~\ref{fig:c2_activity_products})\\
\small FR13 & \small UC13 Update Order Status & \small Main flow of Table~\ref{tab:spec13} & \small ACT-13 (Figure~\ref{fig:c2_activity_admin}), state machine\\
\small FR14 & \small UC14 Find Painting for My Room & \small Specification of UC14 (Chapter~3) & \small ACT-01 (image branch)\\
\bottomrule
\end{longtable}

%% =========================================================================
\section{Consistency Check}\label{sec:2.6}

The models must say the same thing. Lecture~03 takes the requirement ``the customer can cancel an order before
shipment'' and checks whether every analysis artefact shows it correctly. I apply the same check to Họa.
(Table~\ref{tab:consistency}); the last column says where the rule is realised in the later chapters.

\begin{longtable}{@{}P{3.1cm}P{8.0cm}P{3.9cm}@{}}
\caption{Consistency check of the requirement ``cancel before shipment''}\label{tab:consistency}\\
\toprule
\textbf{Artefact} & \textbf{Expected information} & \textbf{Result}\\
\midrule
\endfirsthead
\toprule
\textbf{Artefact} & \textbf{Expected information} & \textbf{Result}\\
\midrule
\endhead
\small Requirement & \small FR11: customers cancel their own order while it is not shipped (PLACED or PROCESSING). & \small Present\\
\small Use case & \small UC11 Cancel Order, extending UC08 View Order. & \small Present (Figure~\ref{fig:usecase})\\
\small Scenario & \small The customer asks to cancel; the system checks ownership, checks the status, sets CANCELLED. & \small Table~\ref{tab:spec11}\\
\small Exception flow & \small An order already shipped is refused. & \small A1 of UC11\\
\small Business rule & \small BR06: cancel only when PLACED or PROCESSING. & \small Table~\ref{tab:rules}\\
\small State machine & \small Transitions PLACED $\rightarrow$ CANCELLED and PROCESSING $\rightarrow$ CANCELLED exist; none from SHIPPED or DELIVERED to CANCELLED. & \small Figure~\ref{fig:stateorder}\\
\small Class model & \small Operation \code{cancel()} of \code{Order} refuses every status except PLACED and PROCESSING; \code{advanceTo(next)} never goes back. & \small Section~4.3\\
\small Test & \small Cancelling a PLACED order succeeds; cancelling the SHIPPED order 20261001 is refused (US08). & \small Chapter~\ref{ch:code}\\
\bottomrule
\end{longtable}

The check found one point that needed a decision: the original requirement said nothing about stock when an order
is cancelled. Instead of leaving the programmer to guess, the point was written down as an open question in the
specification of UC11.

%% =========================================================================
\section{Conclusion}\label{sec:2.10}

Chapter~2 turned the needs of Chapter~\ref{ch:req} into \nFR{} functional and \nNFR{} non-functional requirements
written as ``The system shall\dots'', each with a priority and a source, and into four actors and \nUseCases{} use
cases. It then answered ``when an actor performs a use case, what happens'' for six use cases, each with a main
flow, at least three alternative flows and the diagram that shows its behaviour. Every decision in the activities
comes from a business rule, the trace requirement -- use case -- scenario -- activity is kept for every main
requirement, and the consistency check of ``cancel before shipment'' shows that requirement, use case, scenario,
rule and state machine say the same thing. The domain model, the conceptual and logical data model, the distinction between
\code{CartItem} and \code{OrderItem} and the checkout transformation are the foundation of the design in Chapter~4; the use cases that
involve AI are analysed in Chapter~3.
